\chapter{Composto Mercadológico do BNDES}

\section{Produto}

Com a análise dos produtos de financiamento do BNDES (Banco Nacional de Desenvolvimento Econômico e Social), observam-se condições atrativas e benefícios favoráveis aos tomadores de crédito. Os financiamentos do BNDES podem ser caracterizados como opções de bom valor, pois geralmente oferecem taxas de juros competitivas. Outra vantagem do produto é a possibilidade de prazos mais longos para pagamento, podendo alcançar até cinco anos, incluindo período de carência de até um ano \cite{bndes2026creditodigital}.

\section{Preço}

O Crédito Digital do BNDES apresenta os preços representados pelas taxas de juros de 0,6\% ao ano para micro e pequenas empresas e de 1,25\% ao ano para médias empresas. O valor mínimo do financiamento é de R\$ 1.000,00 por operação de crédito, tornado o produto competitivo em comparação a outros oferecidos no mercado financeiro. Os preços que é observado no valor máximo do financiamento pode chegar a R\$ 10 milhões por cliente final em cada período de 12 meses, observados os seguintes limites por operação:

\begin{itemize}
  \item Micro ou pequena empresa: até R\$ 300 mil;
  \item Média empresa: até R\$ 2,5 milhões \cite{bndes2026creditodigital}.
\end{itemize}

\subsection{Considerações}

Segundo o Painel do Mapa do Cadastro Nacional da Pessoa Jurídica (CNPJ), em julho de 2026 existiam 25.418.820 empresas ativas no Brasil \cite{mdic2026mapacnpj}, demonstrando a relevância do mercado atendido pelo BNDES e a importância de seus produtos e preços financeiros para o desenvolvimento empresarial. O BNDES possui uma ampla linha de produtos e preços financeiros competitivos, atendendo desde microempresas até grandes projetos de infraestrutura. A instituição oferece crédito de longo prazo, com taxas de juros mais baixas e prazos de pagamento mais extensos quando comparados aos praticados por bancos comerciais, contribuindo para o desenvolvimento econômico e social do Brasil \cite{bndes2025estatuto, bndes2026oib}.

\section{Praça (Distribuição)}

A distribuição do BNDES estrutura-se em três caminhos complementares: o apoio indireto, o apoio direto e os canais digitais, que são fortalecidos por ações de capacitação da rede parceira.

\subsection{Apoio indireto: a rede de agentes financeiros}

Para alcançar clientes potenciais em todo o território nacional e no exterior, grande parte das operações do BNDES é realizada de forma indireta, por meio de parceria com uma rede de instituições financeiras credenciadas. Essas instituições, também chamadas de agentes financeiros, são responsáveis pela análise e pela aprovação do financiamento, bem como pela negociação de garantias com o cliente, e assumem o risco de crédito junto ao BNDES \cite{bndes_instituicoes_credenciadas_2026}.

O empresário pode solicitar o financiamento junto à instituição financeira com a qual já possui relacionamento ou escolher outro agente credenciado. Ressalta-se, porém, que cada instituição possui autonomia para decidir se irá oferecer as linhas de crédito do BNDES, de acordo com suas próprias políticas e critérios de concessão. Assim, cabe exclusivamente ao agente financeiro responsável pela operação decidir pela aprovação ou pela rejeição da proposta apresentada \cite{bndes_rede_credenciada_2026}.

Os números confirmam a centralidade desse canal: a rede de agentes financeiros parceiros foi responsável por mais de 69\% do valor aprovado pelo Banco, e 84\% das operações realizadas pela rede parceira foram destinadas a micro, pequenas e médias empresas (MPMEs), o que evidencia o papel da rede na capilaridade das políticas públicas de crédito \cite{bndes_rede_credenciada_2026}.

\subsection{Apoio direto}

Operações de maior porte e recursos não reembolsáveis são distribuídos diretamente pelo Banco. No caso do Fundo de Universalização dos Serviços de Telecomunicações (FUST), por exemplo, o apoio ocorre de três formas: diretamente com o BNDES, por meio de agentes financeiros credenciados ou por meio de fornecedores de equipamentos locais.

Nos programas de natureza social o desenho é semelhante. Em iniciativas como o BNDES Periferias, os recursos não reembolsáveis são disponibilizados e acessados diretamente junto ao BNDES, sem repasse pela rede bancária privada: as instituições interessadas apresentam projetos em chamadas públicas, o próprio Banco analisa e o contrato é firmado diretamente.

\subsection{Canais digitais}

A partir da digitalização, o BNDES passou a encurtar a distância com o cliente final. O Canal MPME é uma plataforma de relacionamento pela internet voltada a pessoas jurídicas com faturamento anual de até R\$ 300 milhões, que permite solicitar financiamento de forma on-line sem a necessidade de visitar previamente um banco repassador. Nela, o empresário identifica as linhas mais adequadas ao seu empreendimento, simula financiamentos, aponta os agentes financeiros de sua preferência e encaminha seu interesse, inclusive por dispositivos móveis. O canal funciona como um hub, no qual são apresentadas as linhas disponíveis e as condições médias de taxa e prazo praticadas por cada instituição, sendo a solicitação concluída no site da instituição escolhida.

Em novembro de 2024, o Banco lançou o BNDES Crédito Digital, que disponibiliza crédito às micro e pequenas empresas por meio dos aplicativos e do internet banking das instituições financeiras parceiras. O empréstimo pode ser contratado em minutos, com o dinheiro na conta do cliente final no mesmo dia, taxa fixa a partir de 1,49\% ao mês e prazo de até 60 meses. O BNDES destinou inicialmente R\$ 1 bilhão ao produto, tendo o Sicredi e o BTG como primeiros agentes a operacionalizá-lo.

\subsection{Apoio ao canal e estrutura física}

A distribuição indireta é sustentada por ações de capacitação: funcionários de instituições parceiras de representação empresarial são capacitados e treinados pelo BNDES para oferecer às empresas informações sobre procedimentos, acesso ao crédito e linhas do Banco \cite{bndes_atendimento_parceiras_2026}. A estrutura física própria, por sua vez, é enxuta, concentrada na sede no Rio de Janeiro e em escritórios como os de Brasília e São Paulo.

\section{Promoção (Comunicação)}

Por ser uma instituição financeira pública, o BNDES adota uma comunicação predominantemente institucional, voltada para o fortalecimento de sua reputação, a promoção da transparência e a construção de um relacionamento sólido com seus diferentes públicos, em vez de priorizar a conversão direta em vendas.

\subsection{Canais de atendimento e relacionamento}

O Banco mantém canais estruturados de contato direto com o público. A Central de Atendimento (0800 702 6337) presta informações sobre produtos, serviços e demais assuntos operacionais e recebe reclamações em primeira instância; a Ouvidoria (0800 702 6307) trata denúncias, sugestões, elogios e reclamações não solucionadas pela Central; e o Serviço de Informação ao Cidadão (0800 887 6000) atende às solicitações amparadas pela Lei de Acesso à Informação. O atendimento ocorre de segunda a sexta-feira, das 8h às 20h. O BNDES ressalta que não possui e-mail de atendimento geral, direcionando os contatos ao formulário do Fale Conosco \cite{bndes_canais_atendimento_2026}.

\subsection{Assessoria de imprensa e mídia própria}

A Agência BNDES de Notícias atua como um canal próprio de comunicação do Banco, reunindo notícias e comunicados institucionais. Os lançamentos de produtos são apresentados em entrevistas coletivas realizadas na sede da instituição, com a participação de diretores do Banco e representantes das instituições parceiras. Essa estratégia contribui para ampliar a divulgação das iniciativas por meio da cobertura da imprensa e da mídia espontânea.

\subsection{Comunicação cooperada com a rede parceira}

Um aspecto distintivo da promoção do BNDES é a comunicação realizada em conjunto com o canal de distribuição. Segundo orientações do próprio \citeonline{bndes_orientacoes_divulgacao_2026}, é comum o cliente buscar informações sobre um financiamento antes mesmo de ir à agência ou de fazer sua solicitação on-line, e os portais da instituição e de seus agentes financeiros figuram entre os principais canais de acesso a esse conteúdo. Por isso, a instituição orienta os parceiros quanto à divulgação de seus produtos, considerando que informações precisas, corretas e alinhadas são fundamentais para fortalecer a parceria e o relacionamento com os clientes \cite{bndes_orientacoes_divulgacao_2026}.

Como procedimento padrão, o BNDES solicita reciprocidade na divulgação de seu endereço eletrônico pelas instituições credenciadas e disponibiliza seu logotipo por meio dos Padrões de Comunicação, garantindo consistência visual da marca ao longo de toda a rede \cite{bndes_orientacoes_divulgacao_2026}.

\subsection{Patrocínio como instrumento de comunicação institucional}

O próprio BNDES considera o patrocínio um instrumento de comunicação institucional voltado ao fortalecimento da marca, à ampliação do diálogo com a sociedade e ao apoio a iniciativas relacionadas ao desenvolvimento econômico, social e ambiental do país \cite{bndes_patrocinios_2026}. Além disso, essa estratégia favorece o relacionamento institucional com organizações públicas e privadas e contribui para a divulgação das iniciativas, produtos e soluções oferecidos pelo Banco \cite{bndes_patrocinios_2026}.

Diante dos projetos culturais, o BNDES apoia iniciativas que contribuam para a valorização de sua marca e que divulguem sua atuação, produtos e serviços junto a públicos de interesse e potenciais clientes, sendo os eventos selecionados mediante chamada pública \cite{bndes_patrocinios_2026}. Em apresentação institucional, o Banco explicitou o encadeamento dessa política: o Planejamento Estratégico orienta o Plano de Comunicação (que define objetivos de comunicação, posicionamento de marca, mensagens-chave, estratégia, ações e monitoramento e mensuração de resultados) o qual, por sua vez, orienta as Diretrizes de Patrocínio. Integra ainda essa frente o Espaço Cultural BNDES, composto por um teatro e uma galeria, com programação inteiramente gratuita à população \cite{bndes_patrocinios_2026}.